% =============================================================================
%  扩散模型基础 —— 学习笔记 / 讲解稿模版
%
%  编译方式（本机 TeX Live 2026）：
%      latexmk -xelatex main.tex          % 一步到位（推荐）
%      make                               % 等价写法
%      xelatex main && biber main && xelatex main && xelatex main
%
%  结构：
%      preamble.tex      全局样式（改外观只动这个文件）
%      sections/*.tex    各章正文，按需增删
%      refs.bib          参考文献
%      figures/          外部图片放这里
% =============================================================================
\documentclass[12pt, a4paper]{ctexart}

% =============================================================================
%  preamble.tex —— 全局样式
%  这里集中放宏包、字体、版式、定理环境、TikZ 样式、代码样式。
%  正文请写在各 sections/*.tex 里；改外观只改这个文件即可。
%  编译： latexmk -xelatex main.tex   （见 Makefile / build.sh）
% =============================================================================


% -----------------------------------------------------------------------------
% 1. 页面与行距
% -----------------------------------------------------------------------------
\usepackage{geometry}
\geometry{
  a4paper,
  left   = 3.2cm, right  = 3.2cm,
  top    = 3.0cm, bottom = 2.8cm,
  headsep = 0.6cm, footskip = 1.0cm
}
\usepackage{setspace}
\setstretch{1.42}          % 中文行距：1.4～1.5 之间比较舒服，可按需调

% 中文段落断行：给 TeX 留一点额外弹性，减少长公式/长词造成的 overfull
\setlength{\emergencystretch}{2.5em}
\tolerance         = 2000
\hbadness          = 10000   % 只关心 overfull，忽略轻微的 underfull
\vbadness          = 10000


% -----------------------------------------------------------------------------
% 2. 字体（XeLaTeX + fontspec）
%    西文用 Latin Modern，与默认的 Computer Modern 数学字体完全统一；
%    中文正文宋体、标题/图表题黑体，与参考笔记的观感一致。
%    若想换成 Times，把 \setmainfont 换成 Times New Roman 即可。
% -----------------------------------------------------------------------------
\usepackage{fontspec}
% 西文衬线族（正文）
\setmainfont{lmroman10-regular.otf}[
  BoldFont       = lmroman10-bold.otf,
  ItalicFont     = lmroman10-italic.otf,
  BoldItalicFont = lmroman10-bolditalic.otf,
  Ligatures      = TeX
]
% 西文无衬线族（标题、图表题）
\setsansfont{Helvetica Neue}
% 西文等宽族（代码）
\setmonofont{Menlo}[Scale=0.92]

% 中文字体
% 注：ctex 已经定义了 \songti \heiti \kaishu \fangsong 等切换命令，
%     这里只需要把四套字体族指向 macOS 上对应的字体即可。
%     Heiti SC 带有真正的 Bold 字面，标题不会出现缺字警告。
\setCJKmainfont{Songti SC}[ItalicFont=STFangsong]   % 宋体；强调用仿宋
\setCJKsansfont{Heiti SC}                          % 黑体（标题）
\setCJKmonofont{STFangsong}                        % 仿宋（等宽场景）


% -----------------------------------------------------------------------------
% 3. 数学
% -----------------------------------------------------------------------------
\usepackage{amsmath}
\usepackage{amssymb}
\usepackage{mathtools}
\usepackage{bm}
\allowdisplaybreaks

% 常用的数学记号（避免正文里反复写长命令）
\providecommand{\E}{\mathbb{E}}                        % 期望
\newcommand{\Ncal}{\mathcal{N}}                        % 高斯分布
\newcommand{\Dkl}{\mathrm{D}_{\mathrm{KL}}}            % KL 散度
\newcommand{\bx}{\bm{x}}
\newcommand{\bz}{\bm{z}}
\newcommand{\beps}{\bm{\varepsilon}}
\newcommand{\bI}{\bm{I}}
\newcommand{\bzero}{\bm{0}}
\newcommand{\bmu}{\bm{\mu}}
\newcommand{\abt}{\bar{\alpha}_t}                      % \bar\alpha_t
\newcommand{\abs}{\bar{\alpha}_s}
\DeclareMathOperator*{\argmax}{arg\,max}
\DeclareMathOperator*{\argmin}{arg\,min}


% -----------------------------------------------------------------------------
% 4. 颜色（先定义，后面 TikZ / 表格 / 代码都要用）
% -----------------------------------------------------------------------------
\usepackage{xcolor}
\definecolor{dmNavy} {HTML}{2B2D5C}   % 深藏青：主线条、标题
\definecolor{dmSteel}{HTML}{3C6E9F}   % 钢蓝：强调、链接
\definecolor{dmLav}  {HTML}{DFDBF4}   % 淡紫：节点底色（对齐参考笔记）
\definecolor{dmSky}  {HTML}{D6E4F7}   % 淡蓝
\definecolor{dmMint} {HTML}{D6EDE7}   % 淡青
\definecolor{dmRose} {HTML}{F7DFDF}   % 淡粉
\definecolor{dmGray} {HTML}{6B6F8A}   % 中性灰


% -----------------------------------------------------------------------------
% 5. 版式细节：标题居中、页码在右上角、图表题样式
% -----------------------------------------------------------------------------
\usepackage{fancyhdr}
\pagestyle{fancy}
\fancyhf{}
\renewcommand{\headrulewidth}{0pt}
\fancyhead[R]{\small\thepage}

% 章节标题：一级居中加粗黑体，二级左对齐
\ctexset{
  section = {
    format     = \sffamily\bfseries\zihao{3}\centering,
    aftername  = \quad,
    beforeskip = 2.0ex plus .3ex minus .2ex,
    afterskip  = 1.4ex plus .2ex,
  },
  subsection = {
    format     = \sffamily\bfseries\zihao{4},
    aftername  = \quad,
    beforeskip = 2.0ex plus .3ex minus .2ex,
    afterskip  = 1.0ex plus .2ex,
  },
  subsubsection = {
    format     = \sffamily\bfseries\zihao{5},
    aftername  = \quad,
    beforeskip = 1.6ex plus .3ex minus .2ex,
    afterskip  = 0.8ex plus .2ex,
  },
}

% 浮动体、图表题
\usepackage{graphicx}
\usepackage{float}
\usepackage{booktabs}
\usepackage{array}
\usepackage{tabularx}
\usepackage{multirow}
\usepackage{caption}
\captionsetup{
  font            = {sf, bf, small},
  labelsep        = colon,
  justification   = centering,
  singlelinecheck = true,
  skip            = 6pt,
}
\captionsetup[table]{position=top}
\newcommand{\tabref}[1]{表~\ref{#1}}
\newcommand{\figref}[1]{图~\ref{#1}}

% 列表间距
\usepackage{enumitem}
\setlist[itemize]{itemsep=0.25ex plus .1ex, topsep=0.6ex plus .2ex, parsep=0pt}
\setlist[enumerate]{itemsep=0.25ex plus .1ex, topsep=0.6ex plus .2ex, parsep=0pt}


% -----------------------------------------------------------------------------
% 6. 定理类环境（正文一律正体，中文不用斜体）
% -----------------------------------------------------------------------------
\usepackage{amsthm}
\newtheoremstyle{cn}
  {1.3ex plus .2ex}% 上方间距
  {1.3ex plus .2ex}% 下方间距
  {\normalfont}%    正文字体：正体
  {}%              缩进
  {\sffamily\bfseries}% 头部字体
  {}%              头部标点
  {0.7em}%         头部后间距
  {}%              头部规格
\theoremstyle{cn}
\newtheorem{definition} {定义}[section]
\newtheorem{proposition}[definition]{命题}
\newtheorem{theorem}    [definition]{定理}
\newtheorem{corollary}  [definition]{推论}
\newtheorem{example}    [definition]{例}
\newtheorem{remark}     [definition]{注}

% 自带的“证明”环境（比 amsthm 的 proof 更适合中文排版）
\newenvironment{pf}[1][证明]{%
  \par\vspace{0.3ex}%
  \noindent{\sffamily\bfseries #1}\quad\ignorespaces
}{%
  \unskip\nobreak\hfill$\square$\par\vspace{1.0ex}%
}


% -----------------------------------------------------------------------------
% 7. TikZ / pgfplots
% -----------------------------------------------------------------------------
\usepackage{tikz}
\usepackage{pgfplots}
\pgfplotsset{compat=1.18}
\usetikzlibrary{
  arrows.meta, positioning, calc, fit, backgrounds,
  shapes.geometric, shapes.misc, decorations.pathreplacing
}

\tikzset{
  % 通用节点：圆角、描边 + 淡紫底（和参考笔记的示意图风格一致）
  dbox/.style  = {draw=dmNavy, line width=0.7pt, fill=dmLav,
                  rounded corners=3pt, inner sep=4pt, minimum height=2.0em,
                  font=\small\sffamily, text=dmNavy, align=center},
  dbox2/.style = {dbox, fill=dmSky},
  dbox3/.style = {dbox, fill=dmMint},
  dbox4/.style = {dbox, fill=dmRose},
  % 打通的箭头
  darr/.style  = {-{Latex[length=2.1mm, width=1.5mm]}, draw=dmNavy, line width=0.7pt},
  dbarr/.style = {{Latex[length=2.1mm, width=1.5mm]}-{Latex[length=2.1mm, width=1.5mm]},
                  draw=dmNavy, line width=0.7pt},
  % 注释文字
  dlab/.style  = {font=\footnotesize\sffamily, text=dmNavy, inner sep=2pt},
  dnum/.style  = {font=\small\itshape, text=dmNavy},
}


% -----------------------------------------------------------------------------
% 8. 算法
% -----------------------------------------------------------------------------
\usepackage{algorithm}
\usepackage{algpseudocode}
\floatname{algorithm}{算法}
\renewcommand{\algorithmicrequire}{\textbf{输入：}}
\renewcommand{\algorithmicensure}{\textbf{输出：}}
\renewcommand{\algorithmicforall}{\textbf{对}}
\renewcommand{\algorithmicdo}{\textbf{do}}
\renewcommand{\algorithmicend}{\textbf{end}}
% 让算法环境的说明文字也有统一的字号
\makeatletter
\renewcommand{\ALG@beginalgorithmic}{\small}
\makeatother


% -----------------------------------------------------------------------------
% 9. 强调框 / 代码
% -----------------------------------------------------------------------------
\usepackage[most]{tcolorbox}
\newtcolorbox{keybox}[1][要点]{
  enhanced, breakable,
  colback  = dmLav!30, colframe = dmNavy,
  boxrule  = 0.7pt, arc = 3pt,
  left = 7pt, right = 7pt, top = 6pt, bottom = 6pt,
  fonttitle = \sffamily\bfseries,
  coltitle  = dmNavy, colbacktitle = dmLav!70,
  title = {#1},
  attach boxed title to top left = {xshift=8pt, yshift=-2pt},
  boxed title style = {boxrule=0pt, arc=2pt},
  top = 12pt,
}

\usepackage{listings}
\lstdefinestyle{pystyle}{
  language        = Python,
  basicstyle      = \ttfamily\footnotesize,
  keywordstyle    = \color{dmSteel}\bfseries,
  commentstyle    = \color{dmGray}\itshape,
  stringstyle     = \color{dmNavy},
  numbers         = left,
  numberstyle     = \tiny\color{dmGray},
  numbersep       = 8pt,
  frame           = single,
  rulecolor       = \color{dmGray!45},
  backgroundcolor = \color{dmGray!6},
  breaklines      = true,
  showstringspaces= false,
  tabsize         = 4,
  xleftmargin     = 1.6em,
  framexleftmargin= 1.6em,
  columns         = fullflexible,
  keepspaces      = true,
  upquote         = true,
}
\lstset{style=pystyle}


% -----------------------------------------------------------------------------
% 10. 超链接与参考文献
%     （hyperref 要放在最后；biblatex 放在 hyperref 之后）
% -----------------------------------------------------------------------------
\usepackage{hyperref}
\hypersetup{
  colorlinks  = true,
  linkcolor   = dmNavy,
  citecolor   = dmSteel,
  urlcolor    = dmSteel,
  bookmarksnumbered = true,
  pdftitle    = {扩散模型基础},
  pdfsubject  = {Diffusion Models, DDPM, 变分推断},
}

\usepackage[
  backend  = biber,
  style    = numeric-comp,
  sortcites= true,
  giveninits = true,
  maxbibnames = 6,
  doi = false, isbn = false,
  url = true
]{biblatex}
\addbibresource{refs.bib}
\setlength{\bibitemsep}{0.4ex}



% -----------------------------------------------------------------------------
%  标题区（想改成 \maketitle 也行，这里手写是为了完全控制行距与字号）
% -----------------------------------------------------------------------------
\newcommand{\doctitle}{扩散模型基础}
\newcommand{\docsubtitle}{从加噪到生成：DDPM 的直观理解、下界推导与采样}
\newcommand{\docdate}{2026 年 9 月}
\newcommand{\docnote}{%
  本文是扩散模型的学习笔记与讲解稿模版。章节的组织方式参考了邱锡鹏老师
  《神经网络与深度学习》的写法——概念先行、直觉在前、公式在后；
  讲解的脉络则参考了李宏毅（Hung-yi Lee）老师的 \emph{DDPM} 系列课件，
  尤其是“\emph{VAE 与扩散模型的对照}”和“\emph{训练时想像中与实际上}”两处。
  文中附上了笔者自己的推导与理解，如有错误或问题，欢迎指正与交流！%
}

\begin{document}

\begin{center}
  {\sffamily\bfseries\zihao{2} \doctitle}\\[0.85em]
  {\sffamily\zihao{5} \docsubtitle}\\[0.55em]
  {\sffamily\zihao{5} \docdate}
\end{center}

\vspace{0.4em}
\docnote

\vspace{0.6em}
\hrule
\vspace{0.6em}

% 想要目录就把下面三行的注释去掉
% \tableofcontents
% \newpage
% \vspace{1em}

\section{生成模型与扩散模型的基本概念}

\subsection{生成模型要解决什么问题}

在监督学习里，我们关心的是从输入到输出的映射：给定样本 $(\bx,y)$，学一个函数 $f$，
让它在新样本上预测得准。这是\emph{判别}的思路，模型刻画的是条件分布 $p(y\mid\bx)$。
生成模型问的是另一个问题：\emph{数据本身}服从什么分布？

记真实数据分布为 $p_{\text{data}}(\bx)$。手头只有从它当中独立抽取的 $N$ 个样本
$\{\bx^{(n)}\}_{n=1}^{N}$（一堆图片、一段语音、一句文本），而 $p_{\text{data}}$ 的具体
形式是未知的。我们希望学到一个模型 $p_\theta(\bx)$，使得\emph{从 $p_\theta$ 采样出来的东西
与训练集里的东西无法区分}。按最大似然估计，
\begin{equation}
  \begin{aligned}
    \theta^\star
    &= \argmax_\theta \prod_{n=1}^{N} p_\theta\bigl(\bx^{(n)}\bigr)
     = \argmax_\theta \sum_{n=1}^{N} \log p_\theta\bigl(\bx^{(n)}\bigr) \\
    &\approx \argmax_\theta\ \E_{\bx\sim p_{\text{data}}}\bigl[\log p_\theta(\bx)\bigr].
  \end{aligned}
  \label{eq:mle}
\end{equation}
把期望写成积分，再减掉一个与 $\theta$ 无关的常数项：
\begin{equation}
  \begin{aligned}
    &\argmax_\theta\ \E_{p_{\text{data}}}\bigl[\log p_\theta(\bx)\bigr]
     = \argmax_\theta \int_{\bx} p_{\text{data}}(\bx)\log p_\theta(\bx)\,\mathrm{d}\bx \\
    &\quad -\underbrace{\int_{\bx} p_{\text{data}}(\bx)\log p_{\text{data}}(\bx)\,\mathrm{d}\bx}_{\text{常数}} \\
    &\quad = \argmin_\theta \int_{\bx} p_{\text{data}}(\bx)
       \log\frac{p_{\text{data}}(\bx)}{p_\theta(\bx)}\,\mathrm{d}\bx \\
    &\quad = \argmin_\theta \Dkl\bigl(p_{\text{data}}\,\|\,p_\theta\bigr).
  \end{aligned}
  \label{eq:mle-kl}
\end{equation}
也就是说，\emph{最大似然等价于最小化真实分布与模型分布之间的 KL 散度}\cite{lee-ddpm}。
这条结论把“训好一个生成模型”还原成了一个统计问题，也是后面所有推导的出发点。

问题在于 $p_\theta(\bx)$ 本身通常算不出来：它往往是某个隐变量积分掉之后的结果，积分不可解。
所有生成模型的差别，就落在“如何绕开这个积分”上。一个统一的写法是：先取一个简单、易采样的
先验分布 $\bz\sim p(\bz)$（通常取标准高斯），再用一个网络 $G_\theta$ 把它变换成数据
\begin{equation}
  \bz \sim \Ncal(\bzero,\bI), \qquad \bx = G_\theta(\bz),
  \label{eq:gen-general}
\end{equation}
差别只在 $G_\theta$ 的结构与训练目标。表~\ref{tab:gen-compare} 给出了几条主要路线的对比。

\begin{table}[htbp]
  \centering
  \caption{几类主流生成模型的对比}
  \label{tab:gen-compare}
  \small
  \renewcommand{\arraystretch}{1.35}
  \setlength{\tabcolsep}{5pt}
  \begin{tabular}{@{}l p{3.0cm} l l p{3.1cm}@{}}
    \toprule
    模型 & 生成方式 & 似然 & 训练 & 特点 \\
    \midrule
    GAN        & 一步映射 $G_\theta(\bz)$            & 不可计算 & 对抗博弈，易不稳定 & 快，但易模式塌缩 \\
    VAE        & 一步映射 $+$ 隐变量                 & 下界     & 稳定               & 快，但结果偏模糊 \\
    流模型     & 可逆映射，$\bz\leftrightarrow\bx$   & 精确     & 稳定               & 快，但受维数与可逆性约束 \\
    扩散模型   & 从噪声出发\emph{迭代去噪} $T$ 次     & 下界     & 稳定               & 慢，但质量高、多样性强 \\
    \bottomrule
  \end{tabular}
\end{table}

\subsection{扩散模型的两个过程}

扩散模型的做法是把式~\eqref{eq:gen-general} 中“一步到位”的映射拆成许多步\cite{sohl2015deep,ho2020ddpm}，
它包含一正一反两个过程（图~\ref{fig:two-process}）：

\begin{itemize}
  \item \textbf{前向过程（加噪）}：预先设定、\emph{无参数}。从干净数据 $\bx_0$ 出发，
        每一步叠加上一小撮高斯噪声，共走 $T$ 步，最终 $\bx_T$ 几乎变成纯噪声，
        分布收敛到 $\Ncal(\bzero,\bI)$；
  \item \textbf{反向过程（去噪）}：\emph{需要学习}。从 $\bx_T\sim\Ncal(\bzero,\bI)$ 出发，
        每一步让网络预测“这一步大概被加进了什么噪声”并把它减掉，逐步还原出 $\bx_0$。
\end{itemize}

\begin{figure}[htbp]
  \centering
  \begin{tikzpicture}[x=1cm, y=1cm]
    % 一排方块：颜色由浅到深，表示噪声逐步吞掉信号
    \foreach \k/\sh in {0/6, 1/20, 3/40, 4/55, 6/80}{%
      \node[draw=dmNavy, line width=0.7pt, rounded corners=2pt,
            minimum size=0.92cm, inner sep=0pt,
            fill=dmNavy!\sh!white] (s\k) at (1.45*\k, 0) {};
    }
    \node[dlab] at (1.45*2, 0) {$\cdots$};
    \node[dlab] at (1.45*5, 0) {$\cdots$};
    % 节点下方的记号
    \node[dlab, below=1pt of s0] {$\bx_0$};
    \node[dlab, below=1pt of s1] {$\bx_1$};
    \node[dlab, below=1pt of s3] {$\bx_{t-1}$};
    \node[dlab, below=1pt of s4] {$\bx_t$};
    \node[dlab, below=1pt of s6] {$\bx_T$};
    \node[dlab, below=1pt of s0, yshift=-0.34cm] {干净数据};
    \node[dlab, below=1pt of s6, yshift=-0.34cm] {纯噪声};
    % 上：前向
    \draw[darr] (-0.46, 1.05) -- (9.16, 1.05);
    \node[dlab, fill=white, inner sep=3pt] at (4.35, 1.05)
      {$q(\bx_t\mid\bx_{t-1})$：逐步加噪（前向过程，无参数）};
    % 下：反向
    \draw[darr] (9.16, -1.25) -- (-0.46, -1.25);
    \node[dlab, fill=white, inner sep=3pt] at (4.35, -1.25)
      {$p_\theta(\bx_{t-1}\mid\bx_t)$：逐步去噪（反向过程，需要学习）};
  \end{tikzpicture}
  \caption{扩散模型的前向过程与反向过程。方块由浅到深表示噪声逐渐增大}
  \label{fig:two-process}
\end{figure}

值得强调的是：前向过程每一步都在丢失信息，它是一个\emph{不可逆}的物理过程，就像一滴墨水滴进清水后
逐渐扩散均匀、熵在增加。我们并不去求它的数学逆，而是让网络学一个\emph{统计意义}上的逆——
给定当前含噪样本，估计这一步被加入的噪声，然后减掉它。这也正是 DDPM 全名的由来：
Denoising Diffusion Probabilistic Models，\emph{去噪扩散概率模型}\cite{ho2020ddpm}。

换一个角度，扩散模型是一种\emph{自回归}式的生成：把“一次到位”改成“$N$ 次到位”\cite{lee-ddpm}。
GAN 与 VAE 一步就把 $\bz$ 变成 $\bx$，扩散模型则用 $T$ 次微小的、\emph{几乎确定}的去噪步骤
完成同一件事。代价是采样变慢（需要 $T$ 次网络前向），换来的是训练的稳定与生成质量的提升。

值得注意的是，这里的“去噪”每一步都只处理一点点噪声，于是每一步的子任务都极其简单，
用一个小网络就能拟合得很好；把 $T$ 步串起来，整体就成了一个表达能力极强的生成器。
这与“把难题拆成一串简单题”的通用思路是一致的。

\subsection{记号约定}

后面会反复出现下列记号，先集中列在这里，读时可随时回查。

\begin{itemize}
  \item $\bx_0\in\mathbb{R}^d$：干净数据（$\bx_0\sim p_{\text{data}}$）；
  \item $\bx_1,\dots,\bx_T\in\mathbb{R}^d$：逐步含噪的隐变量，\emph{与数据同维}（这一点与 VAE 不同）；
  \item $T$：扩散步数，常见 200～1000；
  \item $\beta_t\in(0,1)$：第 $t$ 步的噪声方差；$\{\beta_1,\dots,\beta_T\}$ 称为\emph{噪声调度}（noise schedule）；
  \item $\alpha_t \coloneqq 1-\beta_t$，$\abt\coloneqq\prod_{s=1}^{t}\alpha_s$（注意是\emph{累乘}）；
  \item $\beps\sim\Ncal(\bzero,\bI)$：标准高斯噪声；
  \item $\theta$：网络参数；$\beps_\theta(\bx_t,t)$ 表示噪声预测网络的输出。
\end{itemize}

表~\ref{tab:four-gaussians} 是本笔记的一个“地图”：扩散模型里出现的所有分布都是高斯，
整篇推导其实就是在算它们的均值与方差。

\begin{table}[htbp]
  \centering
  \caption{四个核心分布：全部都是高斯}
  \label{tab:four-gaussians}
  \small
  \renewcommand{\arraystretch}{1.9}
  \setlength{\tabcolsep}{5pt}
  \begin{tabular}{@{}p{4.4cm} p{5.4cm} p{3.0cm}@{}}
    \toprule
    分布 & 均值 & 方差 \\
    \midrule
    前向单步 $q(\bx_t\mid\bx_{t-1})$
      & $\sqrt{\alpha_t}\,\bx_{t-1}$ & $\beta_t\bI$ \\
    前向闭式 $q(\bx_t\mid\bx_0)$
      & $\sqrt{\abt}\,\bx_0$ & $(1-\abt)\bI$ \\
    真实后验 $q(\bx_{t-1}\mid\bx_t,\bx_0)$
      & $\dfrac{\sqrt{\bar\alpha_{t-1}}\,\beta_t\,\bx_0+\sqrt{\alpha_t}\,(1-\bar\alpha_{t-1})\,\bx_t}{1-\abt}$
      & $\dfrac{1-\bar\alpha_{t-1}}{1-\abt}\,\beta_t\bI$ \\
    反向过程 $p_\theta(\bx_{t-1}\mid\bx_t)$
      & $\bmu_\theta(\bx_t,t)$（网络给出） & $\sigma_t^2\bI$（可取 $\beta_t$ 或 $\tilde\beta_t$） \\
    \bottomrule
  \end{tabular}
\end{table}

\section{前向扩散过程}

\subsection{单步加噪}

\begin{definition}[前向扩散过程]
给定噪声调度 $\{\beta_1,\dots,\beta_T\}$，$\beta_t\in(0,1)$，前向过程按如下\emph{马尔可夫链}
逐步给数据加噪：
\begin{equation}
  q(\bx_t\mid\bx_{t-1})
  = \Ncal\Bigl(\bx_t;\ \sqrt{1-\beta_t}\,\bx_{t-1},\ \beta_t\bI\Bigr),
  \qquad t=1,2,\dots,T,
  \label{eq:forward-step}
\end{equation}
等价地写成重参数化的形式
\begin{equation}
  \bx_t = \sqrt{1-\beta_t}\,\bx_{t-1} + \sqrt{\beta_t}\,\beps_{t-1},
  \qquad \beps_{t-1}\sim\Ncal(\bzero,\bI)\ \text{相互独立}.
  \label{eq:forward-step-reparam}
\end{equation}
记 $\alpha_t = 1-\beta_t$，则整条链的联合分布为
$q(\bx_{1:T}\mid\bx_0)=\prod_{t=1}^{T}q(\bx_t\mid\bx_{t-1})$。
\end{definition}

\begin{remark}[为什么系数是 $\sqrt{1-\beta_t}$ 而不是 $1$]
式~\eqref{eq:forward-step-reparam} 里信号的系数与噪声的系数是一对平方和为 $1$ 的数：
$(\sqrt{1-\beta_t})^2+(\sqrt{\beta_t})^2=1$。若 $\bx_{t-1}$ 的各维独立、均值为 $0$、方差为 $1$，
那么 $\bx_t$ 的方差为 $(1-\beta_t)\cdot 1+\beta_t=1$，\emph{方差始终为 $1$}。
这个约定称为保方差（variance preserving）。反之若写成 $\bx_t=\bx_{t-1}+\sqrt{\beta_t}\beps$，
方差会随 $t$ 线性增长，到 $T=1000$ 时信号会被彻底淹没，网络也就学不动了。
\end{remark}

\subsection{任意时刻的闭式解}

逐层递推 $T$ 次显然不可行。但幸运的是，由于每一步都是高斯，整条链可以被\textbf{一步折叠}，
直接写出任意 $t$ 时刻 $\bx_t$ 关于 $\bx_0$ 的分布——这是扩散模型最核心的一条性质。

\begin{proposition}[前向过程的闭式解]
\label{prop:forward-closed}
对任意 $t=1,\dots,T$，
\begin{equation}
  q(\bx_t\mid\bx_0)=\Ncal\Bigl(\bx_t;\ \sqrt{\abt}\,\bx_0,\ (1-\abt)\bI\Bigr),
  \label{eq:forward-closed}
\end{equation}
即
\begin{equation}
  \boxed{\ \bx_t=\sqrt{\abt}\,\bx_0+\sqrt{1-\abt}\,\beps,\qquad \beps\sim\Ncal(\bzero,\bI).\ }
  \label{eq:forward-closed-reparam}
\end{equation}
\end{proposition}

\begin{pf}
对 $t$ 归纳。$t=1$ 时即为定义式~\eqref{eq:forward-step-reparam}。
设 $t-1$ 时命题成立，即 $\bx_{t-1}=\sqrt{\bar\alpha_{t-1}}\,\bx_0+\sqrt{1-\bar\alpha_{t-1}}\,\beps'$，
代入单步公式并合并同类项：
\begin{equation}
  \begin{aligned}
    \bx_t
    &= \sqrt{\alpha_t}\,\bx_{t-1}+\sqrt{\beta_t}\,\beps''
     = \sqrt{\alpha_t\bar\alpha_{t-1}}\,\bx_0
       + \underbrace{\sqrt{\alpha_t(1-\bar\alpha_{t-1})}\,\beps'+\sqrt{\beta_t}\,\beps''}_{\text{两个独立零均值高斯之和}} .
  \end{aligned}
  \label{eq:forward-induction}
\end{equation}
由附录~\ref{app:gauss-sum}，两个独立的零均值高斯之和仍是零均值高斯，其方差为两者方差之和：
\begin{equation}
  \alpha_t\,(1-\bar\alpha_{t-1})+\beta_t
  = \alpha_t-\alpha_t\bar\alpha_{t-1}+1-\alpha_t
  = 1-\alpha_t\bar\alpha_{t-1}
  = 1-\abt .
  \label{eq:forward-variance}
\end{equation}
又 $\sqrt{\alpha_t\bar\alpha_{t-1}}=\sqrt{\abt}$，故 $\bx_t$ 服从式~\eqref{eq:forward-closed}。
归纳完成。
\end{pf}

\begin{corollary}[起点的分布是已知的]
令 $t=T$ 得 $q(\bx_T\mid\bx_0)=\Ncal(\bx_T;\sqrt{\bar\alpha_T}\,\bx_0,(1-\bar\alpha_T)\bI)$。
只要调度使得 $\bar\alpha_T\to 0$，就有 $\bx_T\to\Ncal(\bzero,\bI)$ 且\emph{与 $\bx_0$ 无关}。
这保证了反向过程有一个已知的、可以直接采样的起点：从标准高斯里取一个噪声，就能开始生成。
\end{corollary}

\subsection{训练时的「想像中」与「实际上」}

按定义，要得到一个含噪样本 $\bx_t$ 似乎必须老老实实走 $t$ 步。但命题~\ref{prop:forward-closed} 告诉我们
可以\emph{一步算出}。这一点在教学材料中常被专门拿出来强调，因为它是扩散模型能够高效训练的
关键，也是初学者最容易忽略的地方\cite{lee-ddpm}。

\begin{keybox}[一句话记住 $\bx_t$ 的来源]
\begin{center}
  \textbf{想像中}：$\bx_0 \xrightarrow{\ \sqrt{\beta_1}\beps_1\ } \bx_1
                  \xrightarrow{\ \sqrt{\beta_2}\beps_2\ } \bx_2
                  \xrightarrow{\ \cdots\ } \bx_t$
                  \quad（逐步加噪，$t$ 次）\\[0.9ex]
  \textbf{实际上}：$\bx_0,\ \beps \xrightarrow{\ \text{式}\eqref{eq:forward-closed-reparam}\ }
                  \bx_t$ \quad（一次算出，$O(1)$）
\end{center}
训练时随机采一个 $t$、再采一个 $\beps$，直接用闭式解造出 $\bx_t$，
并以当初加进去的 $\beps$ 作为监督目标。全程只做一次加噪，与 $t$ 无关。
\end{keybox}

\subsection{噪声调度与信噪比}

$\abt$ 随 $t$ 单调递减：$t$ 越大，剩下的信号越少。刻画“信号还剩多少”的常用量是\emph{信噪比}
\begin{equation}
  \mathrm{SNR}(t)=\frac{\abt}{1-\abt},
  \label{eq:snr}
\end{equation}
它同样单调递减，从 $t=1$ 时的约 $10^4$ 一路降到 $t=1000$ 时的 $10^{-5}$ 量级。
$\{\beta_t\}$ 的取法称为噪声调度，常见两种\cite{ho2020ddpm,nichol2021improved}：

\begin{itemize}
  \item \textbf{线性调度}：$\beta_t$ 从 $\beta_1=10^{-4}$ 线性增到 $\beta_T=2\times10^{-2}$，$T=1000$。
        实现最简单，但两端变化过快、中间过于平坦；
  \item \textbf{余弦调度}：令
        $\abt=\dfrac{\cos^2\!\bigl(\frac{\pi}{2}\cdot\frac{t/T+s}{1+s}\bigr)}{\cos^2\!\bigl(\frac{\pi}{2}\cdot\frac{s}{1+s}\bigr)}$，
        $s\approx 0.008$。它让 $\abt$ 的下降更均匀，低噪声阶段保留信号更久，实践中生成质量更好。
\end{itemize}

\begin{figure}[htbp]
  \centering
  \begin{minipage}[c]{0.47\textwidth}
    \centering
    \resizebox{\linewidth}{!}{%
    \begin{tikzpicture}
      \begin{axis}[
        width=\textwidth, height=5.2cm,
        xlabel={扩散步 $t$}, ylabel={$\abt$},
        xmin=0, xmax=1000, ymin=0, ymax=1.06,
        domain=0:1000, samples=161,
        axis lines=left,
        every axis label/.style={font=\footnotesize\sffamily, text=dmNavy},
        tick label style={font=\scriptsize, text=dmGray},
        legend style={font=\scriptsize, draw=dmGray!40, fill=white,
                      at={(0.97,0.97)}, anchor=north east},
        label style={font=\scriptsize},
      ]
        \addplot[dmNavy, very thick]
          {exp(-(x*0.0001 + 0.0199*x*(x-1)/1998))};
        \addlegendentry{线性调度}
        \addplot[dmSteel, very thick, dashed]
          {(cos(deg(pi/2*(x/1000+0.008)/1.008)))^2/(cos(deg(pi/2*0.008/1.008)))^2};
        \addlegendentry{余弦调度}
      \end{axis}
    \end{tikzpicture}}
    \\[-0.4em]
    {\footnotesize\sffamily 信号保留比例 $\abt$}
  \end{minipage}\hfill
  \begin{minipage}[c]{0.47\textwidth}
    \centering
    \resizebox{\linewidth}{!}{%
    \begin{tikzpicture}
      \begin{axis}[
        width=\textwidth, height=5.2cm,
        xlabel={扩散步 $t$}, ylabel={$\mathrm{SNR}(t)$},
        xmin=1, xmax=1000, ymode=log,
        ymin=1e-5, ymax=1e5,
        domain=1:1000, samples=161,
        axis lines=left,
        every axis label/.style={font=\footnotesize\sffamily, text=dmNavy},
        tick label style={font=\scriptsize, text=dmGray},
        legend style={font=\scriptsize, draw=dmGray!40, fill=white,
                      at={(0.97,0.03)}, anchor=south east},
        label style={font=\scriptsize},
      ]
        \addplot[dmNavy, very thick]
          {exp(-(x*0.0001 + 0.0199*x*(x-1)/1998))
           /(1-exp(-(x*0.0001 + 0.0199*x*(x-1)/1998)))};
        \addlegendentry{线性调度}
        \addplot[dmSteel, very thick, dashed]
          {((cos(deg(pi/2*(x/1000+0.008)/1.008)))^2/(cos(deg(pi/2*0.008/1.008)))^2)
           /(1-((cos(deg(pi/2*(x/1000+0.008)/1.008)))^2/(cos(deg(pi/2*0.008/1.008)))^2))};
        \addlegendentry{余弦调度}
      \end{axis}
    \end{tikzpicture}}
    \\[-0.4em]
    {\footnotesize\sffamily 信噪比 $\mathrm{SNR}=\abt/(1-\abt)$}
  \end{minipage}
  \caption{两种噪声调度下 $\abt$ 与信噪比随扩散步的变化。线性调度的曲线由
    $\abt\approx\exp\bigl(-\sum_{s\le t}\beta_s\bigr)$ 给出}
  \label{fig:schedule}
\end{figure}

图~\ref{fig:schedule} 对比了两种调度。可以看到线性调度在 $t\approx 250$ 时就已经把一半以上的信号
换成了噪声，而余弦调度要到 $t\approx 500$ 才降到同一水平。这意味着余弦调度把更多的“容量”
留给了中等噪声区间，而这个区间恰恰是反向过程最难学、也最影响生成细节的部分。

\section{从 VAE 的似然下界到 DDPM 的下界}

\subsection{似然算不出来，那就找一个下界}

回到式~\eqref{eq:gen-general}。以 VAE 为例，模型分布是
\begin{equation}
  p_\theta(\bx)=\int_{\bz} p(\bz)\,p_\theta(\bx\mid\bz)\,\mathrm{d}\bz,
  \qquad
  p_\theta(\bx\mid\bz)\propto\exp\!\Bigl(-\tfrac{1}{2\sigma^2}\bigl\|G_\theta(\bz)-\bx\bigr\|^2\Bigr),
  \label{eq:vae-likelihood}
\end{equation}
即解码器输出的是一个均值为 $G_\theta(\bz)$ 的高斯。这个积分一般没有解析解，也不便于直接最大化。
标准的出路是引入一个\emph{任意}的辅助分布 $q(\bz\mid\bx)$（在 VAE 里它就是编码器），
再用 Jensen 不等式把对数挪到期望里面去：
\begin{equation}
  \begin{aligned}
    \log p_\theta(\bx)
    &= \log \int_{\bz} p_\theta(\bx,\bz)\,\mathrm{d}\bz
     = \log \int_{\bz} q(\bz\mid\bx)\,\frac{p_\theta(\bx,\bz)}{q(\bz\mid\bx)}\,\mathrm{d}\bz \\
    &\ge \int_{\bz} q(\bz\mid\bx)\log\frac{p_\theta(\bx,\bz)}{q(\bz\mid\bx)}\,\mathrm{d}\bz
     = \E_{q(\bz\mid\bx)}\Bigl[\log\frac{p_\theta(\bx,\bz)}{q(\bz\mid\bx)}\Bigr]
     \eqqcolon \mathcal{L}(\bx),
  \end{aligned}
  \label{eq:elbo}
\end{equation}
其中第一步用了 $\log\int=\log\E$，第二步用了 Jensen 不等式 $\log\E[\cdot]\ge\E[\log\cdot]$。
这一项称为\emph{证据下界}（evidence lower bound, ELBO），取等号当且仅当
$q(\bz\mid\bx)=p_\theta(\bz\mid\bx)$。把 $\log\frac{p_\theta(\bx,\bz)}{q(\bz\mid\bx)}$ 拆开，得到
更熟悉的形式
\begin{equation}
  \mathcal{L}(\bx)
  = \underbrace{\E_{q(\bz\mid\bx)}\bigl[\log p_\theta(\bx\mid\bz)\bigr]}_{\text{重构项}}
  - \underbrace{\Dkl\bigl(q(\bz\mid\bx)\,\|\,p(\bz)\bigr)}_{\text{把编码器拉向先验}}.
  \label{eq:elbo-vae}
\end{equation}

\subsection{扩散模型的似然}

扩散模型把“一个隐变量 $\bz$”换成“一串隐变量 $\bx_1,\dots,\bx_T$”，于是
\begin{equation}
  \begin{aligned}
    p_\theta(\bx_0)
    &=\int p(\bx_T)\prod_{t=1}^{T}p_\theta(\bx_{t-1}\mid\bx_t)\,\mathrm{d}\bx_{1:T}, \\
    p_\theta(\bx_{t-1}\mid\bx_t)
    &\propto\exp\!\Bigl(-\tfrac{1}{2\sigma_t^2}
      \bigl\|G_\theta(\bx_t)-\bx_{t-1}\bigr\|^2\Bigr),
  \end{aligned}
  \label{eq:ddpm-likelihood}
\end{equation}
形式上与式~\eqref{eq:vae-likelihood} 一模一样，只是积分从一个隐变量变成了一串。
沿着式~\eqref{eq:elbo} 的路子照抄，把辅助分布取成\emph{前向过程}
$q(\bx_{1:T}\mid\bx_0)=\prod_{t=1}^{T}q(\bx_t\mid\bx_{t-1})$，就得到

\begin{equation}
  \log p_\theta(\bx_0)\ \ge\
  \mathcal{L}(\bx_0)=\E_{q(\bx_{1:T}\mid\bx_0)}
  \Bigl[\log\frac{p_\theta(\bx_{0:T})}{q(\bx_{1:T}\mid\bx_0)}\Bigr].
  \label{eq:ddpm-elbo}
\end{equation}

这里有一个非常漂亮的类比\cite{lee-ddpm}：\textbf{前向过程就是 VAE 里的编码器}。
区别在于，这个“编码器”是人为设定、逐层加噪的，\emph{不含任何参数、也不需要学习}；
而“解码器”$p_\theta(\bx_{t-1}\mid\bx_t)$ 需要学习，且不是一层，而是 $T$ 层。
表~\ref{tab:vae-vs-diff} 把两者的对应关系列了出来。

\begin{table}[htbp]
  \centering
  \caption{VAE 与扩散模型的逐项对照}
  \label{tab:vae-vs-diff}
  \small
  \renewcommand{\arraystretch}{1.35}
  \setlength{\tabcolsep}{6pt}
  \begin{tabular}{@{}l p{4.6cm} p{5.6cm}@{}}
    \toprule
     & VAE & 扩散模型 \\
    \midrule
    隐变量       & $\bz$，低维           & $\bx_1,\dots,\bx_T$，共 $T$ 个，与数据同维 \\
    编码器       & $q_\phi(\bz\mid\bx)$，需要学习 & $q(\bx_{1:T}\mid\bx_0)$，固定、无参数 \\
    解码器       & $p_\theta(\bx\mid\bz)$，一层     & $p_\theta(\bx_{t-1}\mid\bx_t)$，$T$ 层，需学习 \\
    训练目标     & 式~\eqref{eq:elbo-vae} 的 ELBO & 式~\eqref{eq:ddpm-elbo} 的 ELBO \\
    主要代价     & 结果偏模糊、后验易坍缩 & 采样需要 $T$ 次网络前向，慢 \\
    \bottomrule
  \end{tabular}
\end{table}

\subsection{下界的三项分解}

把 $p_\theta(\bx_{0:T})=p(\bx_T)\prod_{t=1}^{T}p_\theta(\bx_{t-1}\mid\bx_t)$ 与
$q(\bx_{1:T}\mid\bx_0)=\prod_{t=1}^{T}q(\bx_t\mid\bx_{t-1})$ 代入式~\eqref{eq:ddpm-elbo}，
再对 $t\ge 2$ 的每一对相邻项用贝叶斯公式
\begin{equation}
  q(\bx_t\mid\bx_{t-1})
  = \frac{q(\bx_{t-1}\mid\bx_t,\bx_0)\,q(\bx_t\mid\bx_0)}{q(\bx_{t-1}\mid\bx_0)}
  \qquad(\text{因为 } q(\bx_t\mid\bx_{t-1},\bx_0)=q(\bx_t\mid\bx_{t-1})),
  \label{eq:bayes-trick}
\end{equation}
$q(\bx_t\mid\bx_0)/q(\bx_{t-1}\mid\bx_0)$ 这两项在求和时会\emph{相邻相消}（telescoping），
最后只剩首尾两项。整理后得到 DDPM 论文中的经典分解\cite{ho2020ddpm}（完整步骤见附录~\ref{app:elbo}）：
\begin{equation}
  \log p_\theta(\bx_0)\ \ge\ -\bigl(L_T+\textstyle\sum_{t=2}^{T}L_{t-1}+L_0\bigr),
  \label{eq:elbo-three}
\end{equation}
其中三项分别是
\begin{equation}
  \left\{
  \begin{aligned}
    L_T      &= \Dkl\bigl(q(\bx_T\mid\bx_0)\,\|\,p(\bx_T)\bigr),
      &&\text{常数项，与 }\theta\text{ 无关}; \\
    L_{t-1}  &= \Dkl\bigl(q(\bx_{t-1}\mid\bx_t,\bx_0)\,\|\,p_\theta(\bx_{t-1}\mid\bx_t)\bigr),
      &&\text{主项，共 }T-1\text{ 项}; \\
    L_0      &= -\log p_\theta(\bx_0\mid\bx_1),
      &&\text{重构项}.
  \end{aligned}
  \right.
  \label{eq:elbo-terms}
\end{equation}

于是\textbf{训练扩散模型 = 最小化 $L_T+\sum_{t=2}^{T}L_{t-1}+L_0$}。
逐项看一眼：

\begin{itemize}
  \item $L_T$ 只由噪声调度决定，$p(\bx_T)=\Ncal(\bzero,\bI)$ 是固定的先验，
        只要调度让 $\bar\alpha_T\to 0$，$q(\bx_T\mid\bx_0)$ 就已经接近 $\Ncal(\bzero,\bI)$，
        这一项几乎为 $0$，所以直接丢掉；
  \item $L_{t-1}$ 是全部的工作量所在。它在“网络给出的反向分布”与“真实后验”之间算 KL，
        而\emph{真实后验 $q(\bx_{t-1}\mid\bx_t,\bx_0)$ 是一个已知的、可以解析写出的高斯}——
        这正是下一节要用的关键事实；
  \item $L_0$ 是最后一跳的似然，实践中通常把它并入 $t=1$ 的那一项，
        或对像素做离散化后按 $\{0,\dots,255\}$ 的分段高斯处理。
\end{itemize}

\begin{keybox}[本节的落点]
$L_{t-1}$ 里的两个分布都是高斯，且方差都可以固定为常数，于是 KL 散度\emph{退化成两个均值之间的
平方误差}。所以“用变分推断推出的下界”最终会变成一个极其朴素的回归问题——
这既是扩散模型好训的原因，也是它一开始让人意外的地方。
\end{keybox}

\section{反向过程与训练目标}

\subsection{反向过程的参数化}

反向过程是真正需要学习的部分。它的形式是

\begin{equation}
  p_\theta(\bx_{0:T})=p(\bx_T)\prod_{t=1}^{T}p_\theta(\bx_{t-1}\mid\bx_t),
  \qquad
  p(\bx_T)=\Ncal(\bzero,\bI),
  \label{eq:reverse-chain}
\end{equation}
其中每一步取\emph{均值由网络给出、方差固定}的高斯：
\begin{equation}
  p_\theta(\bx_{t-1}\mid\bx_t)
  = \Ncal\Bigl(\bx_{t-1};\ \bmu_\theta(\bx_t,t),\ \sigma_t^2\bI\Bigr),
  \qquad t=T,T-1,\dots,1 .
  \label{eq:reverse-step}
\end{equation}
只让网络决定均值、方差取一个预先定好的 $\sigma_t^2$，是 DDPM 的关键简化：一方面让后面
$L_{t-1}$ 中的两个高斯方差相同，KL 有干净的闭式；另一方面少了一个学习目标，训练更稳。

\begin{remark}[和 VAE 解码器的差别]
VAE 的解码器把低维的 $\bz$ 映回高维的 $\bx$，一次跳很远，所以容易模糊；
扩散模型的每一步只把 $\bx_t$ 挪动一点点到 $\bx_{t-1}$，输入输出几乎逐像素对应，
于是网络实际上是在做“去一点点噪声”这件非常简单的事。$T$ 步串起来，累计的效果就很强。
\end{remark}

\subsection{真实后验的闭式解}

上一节说过，$L_{t-1}$ 里出现的 $q(\bx_{t-1}\mid\bx_t,\bx_0)$ 是已知的。
下面把它写出来——这是整篇推导里最需要动手的一块代数。

\begin{proposition}[前向过程的后验]
\label{prop:posterior}
对 $t\ge 2$，
\begin{equation}
  q(\bx_{t-1}\mid\bx_t,\bx_0)
  = \Ncal\Bigl(\bx_{t-1};\ \tilde{\bmu}_t(\bx_t,\bx_0),\ \tilde\beta_t\bI\Bigr),
  \label{eq:posterior}
\end{equation}
其中
\begin{equation}
  \tilde{\bmu}_t(\bx_t,\bx_0)
  = \frac{\sqrt{\bar\alpha_{t-1}}\,\beta_t\,\bx_0+\sqrt{\alpha_t}\,(1-\bar\alpha_{t-1})\,\bx_t}{1-\abt},
  \qquad
  \tilde\beta_t=\frac{1-\bar\alpha_{t-1}}{1-\abt}\,\beta_t .
  \label{eq:posterior-mean-var}
\end{equation}
\end{proposition}

\begin{pf}
由贝叶斯公式并把 $q(\bx_t\mid\bx_{t-1},\bx_0)=q(\bx_t\mid\bx_{t-1})$ 代入，
\begin{equation}
  q(\bx_{t-1}\mid\bx_t,\bx_0)
  =\frac{q(\bx_t\mid\bx_{t-1})\,q(\bx_{t-1}\mid\bx_0)}{q(\bx_t\mid\bx_0)} ,
  \label{eq:posterior-bayes}
\end{equation}
分母与 $\bx_{t-1}$ 无关，只需看分子。分子的两个因子都是已知高斯：
\begin{equation}
  \begin{aligned}
    q(\bx_t\mid\bx_{t-1})
    &\propto \exp\Bigl(-\frac{\bigl\|\bx_t-\sqrt{\alpha_t}\,\bx_{t-1}\bigr\|^2}{2\beta_t}\Bigr), \\
    q(\bx_{t-1}\mid\bx_0)
    &\propto \exp\Bigl(-\frac{\bigl\|\bx_{t-1}-\sqrt{\bar\alpha_{t-1}}\,\bx_0\bigr\|^2}{2(1-\bar\alpha_{t-1})}\Bigr).
  \end{aligned}
  \label{eq:posterior-factors}
\end{equation}
把指数相加，写成关于 $\bx_{t-1}$ 的二次型
$-\frac{1}{2}\bigl(A\|\bx_{t-1}\|^2-2\,\bx_{t-1}^{\top}\bm{b}\bigr)+C$，其中
\begin{equation}
  \begin{aligned}
    A&=\frac{\alpha_t}{\beta_t}+\frac{1}{1-\bar\alpha_{t-1}}
     =\frac{\alpha_t(1-\bar\alpha_{t-1})+\beta_t}{\beta_t(1-\bar\alpha_{t-1})}
     =\frac{1-\abt}{\beta_t(1-\bar\alpha_{t-1})}, \\
    \bm{b}&=\frac{\sqrt{\alpha_t}}{\beta_t}\,\bx_t
     +\frac{\sqrt{\bar\alpha_{t-1}}}{1-\bar\alpha_{t-1}}\,\bx_0 .
  \end{aligned}
  \label{eq:posterior-ab}
\end{equation}
二次型配方即得高斯 $\Ncal\bigl(\bm{b}/A,\ A^{-1}\bI\bigr)$。代回
$A^{-1}=\dfrac{\beta_t(1-\bar\alpha_{t-1})}{1-\abt}=\tilde\beta_t$，
并整理 $\bm{b}/A$ 便得到式~\eqref{eq:posterior-mean-var}。
\end{pf}

这个命题的意义在于：它把一条 $T$ 步的链，压缩成了\emph{任意一对相邻状态之间的条件分布}。
有了它，$L_{t-1}$ 里“真实后验 vs 网络输出”的对比就变成了两个显式高斯之间的对比。

\subsection{从“预测均值”到“预测噪声”}

有了命题~\ref{prop:posterior}，就可以把 $L_{t-1}$ 写开了。取 $\sigma_t^2=\beta_t$，
由附录~\ref{app:kl-gauss} 的高斯 KL 公式（两者方差相同，只有均值不同）：
\begin{equation}
  L_{t-1}
  =\E_{q(\bx_t\mid\bx_0)}\Bigl[\Dkl\bigl(\Ncal(\tilde{\bmu}_t,\beta_t\bI)\,\big\|\,\Ncal(\bmu_\theta,\beta_t\bI)\bigr)\Bigr]
  =\E_{q(\bx_t\mid\bx_0)}\Bigl[\frac{1}{2\beta_t}\bigl\|\tilde{\bmu}_t-\bmu_\theta(\bx_t,t)\bigr\|^2\Bigr].
  \label{eq:Lt-kl}
\end{equation}
于是最小化 $L_{t-1}$，就是让网络输出的均值去逼近 $\tilde{\bmu}_t$。

问题是 $\tilde{\bmu}_t$ 里含 $\bx_0$，而 $\bx_0$ 恰好是我们想生成的东西，采样时根本拿不到。
出路是利用重参数化把 $\bx_0$ \emph{解出来}：
\begin{equation}
  \bx_t=\sqrt{\abt}\,\bx_0+\sqrt{1-\abt}\,\beps
  \quad\Longrightarrow\quad
  \bx_0=\frac{\bx_t-\sqrt{1-\abt}\,\beps}{\sqrt{\abt}} .
  \label{eq:x0-from-eps}
\end{equation}
代回式~\eqref{eq:posterior-mean-var} 中的 $\tilde{\bmu}_t$ 并整理，得到一个只含 $\bx_t$ 与 $\beps$ 的
漂亮形式
\begin{equation}
  \boxed{\ \tilde{\bmu}_t=\frac{1}{\sqrt{\alpha_t}}
    \Bigl(\bx_t-\frac{\beta_t}{\sqrt{1-\abt}}\,\beps\Bigr).\ }
  \label{eq:mu-from-eps}
\end{equation}
换言之：\emph{“这一步该往哪走”完全由“这一步被加了什么噪声”决定}，不需要知道 $\bx_0$。
既然 $\beps$ 未知，就干脆让网络直接去预测它。取
\begin{equation}
  \bmu_\theta(\bx_t,t)=\frac{1}{\sqrt{\alpha_t}}
  \Bigl(\bx_t-\frac{\beta_t}{\sqrt{1-\abt}}\,\beps_\theta(\bx_t,t)\Bigr),
  \label{eq:mu-theta}
\end{equation}
代回式~\eqref{eq:Lt-kl}，两个均值之差变成
\begin{equation}
  \bigl\|\tilde{\bmu}_t-\bmu_\theta\bigr\|^2
  =\frac{\beta_t^2}{\alpha_t\,(1-\abt)}\,\bigl\|\beps-\beps_\theta(\bx_t,t)\bigr\|^2 ,
  \label{eq:mean-diff}
\end{equation}
于是
\begin{equation}
  L_{t-1}=w_t\,\E\Bigl[\bigl\|\beps-\beps_\theta(\bx_t,t)\bigr\|^2\Bigr],
  \qquad
  w_t=\frac{\beta_t^2}{2\sigma_t^2\,\alpha_t\,(1-\abt)} .
  \label{eq:Lt-weighted}
\end{equation}

\begin{remark}[式~\eqref{eq:mu-from-eps} 的验算]
把 $\beps=(\bx_t-\sqrt{\abt}\bx_0)/\sqrt{1-\abt}$ 代进去：
\begin{equation}
  \frac{1}{\sqrt{\alpha_t}}\Bigl(\bx_t-\frac{\beta_t}{1-\abt}\bigl(\bx_t-\sqrt{\abt}\bx_0\bigr)\Bigr)
  =\frac{\sqrt{\alpha_t}\,(1-\bar\alpha_{t-1})}{1-\abt}\,\bx_t
   +\frac{\sqrt{\bar\alpha_{t-1}}\,\beta_t}{1-\abt}\,\bx_0 ,
\end{equation}
最后恰好是式~\eqref{eq:posterior-mean-var} 中的 $\tilde{\bmu}_t$。
\end{remark}

\subsection{简化损失函数}

最后一步：丢掉权重 $w_t$，并对所有 $t$ 取简单平均。DDPM 论文指出，
去掉 $w_t$ 相当于给不同噪声等级的下界项重新加权，
实验上\emph{反而}能得到更好的生成质量\cite{ho2020ddpm}。这样就得到了扩散模型最著名的那个目标函数：

\begin{equation}
  \boxed{\
  \mathcal{L}_{\text{simple}}(\theta)
  =\E_{\substack{t\sim\mathcal{U}\{1,\dots,T\}\\ \bx_0\sim p_{\text{data}},\ \beps\sim\Ncal(\bzero,\bI)}}
  \Bigl[\bigl\|\beps-\beps_\theta\bigl(\sqrt{\abt}\,\bx_0+\sqrt{1-\abt}\,\beps,\ t\bigr)\bigr\|^2\Bigr].
  \ }
  \label{eq:l-simple}
\end{equation}

读一遍这个式子，它的内容朴素得有点出人意料：

\begin{enumerate}
  \item 取一张干净数据 $\bx_0$；
  \item 随机采一个噪声等级 $t$，再采一撮标准高斯噪声 $\beps$；
  \item 按闭式解合成含噪图 $\bx_t=\sqrt{\abt}\bx_0+\sqrt{1-\abt}\beps$；
  \item 把 $\bx_t$ 和 $t$ 喂给网络，让它猜出刚才那撮噪声 $\beps$；
  \item 用平方误差回传梯度。
\end{enumerate}

一个由变分推断严格推出的下界，最后化简成了“\emph{看图猜噪声}”这样一个回归任务。
它形式上就是一个去噪自编码器（denoising autoencoder）的目标，也与
去噪得分匹配（denoising score matching）等价\cite{vincent2011connection,song2021sde}。

\begin{keybox}[预测噪声 $=$ 估计得分]
由 $q(\bx_t\mid\bx_0)$ 的高斯形式可以直接求对数密度的梯度：
\begin{equation*}
  \nabla_{\bx_t}\log q(\bx_t\mid\bx_0)
  =-\frac{\bx_t-\sqrt{\abt}\,\bx_0}{1-\abt}
  =-\frac{\beps}{\sqrt{1-\abt}} .
\end{equation*}
边际得分 $\nabla_{\bx_t}\log q(\bx_t)=\E\bigl[\nabla_{\bx_t}\log q(\bx_t\mid\bx_0)\bigr]$，
所以“预测噪声 $\beps_\theta$”与“估计数据分布的对数密度梯度”是同一件事的两种说法。
这也解释了为什么扩散模型可以写成随机微分方程、并用概率流 ODE 来采样。
\end{keybox}

\subsection{三种常用的参数化}

式~\eqref{eq:x0-from-eps} 说明 $\bx_0$、$\bx_t$、$\beps$ 三个量之间两两可以互推，
所以网络输出哪一个目标都可以，差别在于不同噪声等级下的数值稳定性。常见三种：

\begin{table}[htbp]
  \centering
  \caption{三种参数化方式及其换算关系}
  \label{tab:param}
  \small
  \renewcommand{\arraystretch}{1.55}
  \setlength{\tabcolsep}{4pt}
  \begin{tabular}{@{}l l p{4.9cm} p{3.1cm}@{}}
    \toprule
    参数化 & 网络输出 & 反解 & 特点 \\
    \midrule
    $\beps$-参数化 & $\beps$
      & $\bx_0=\dfrac{\bx_t-\sqrt{1-\abt}\,\beps}{\sqrt{\abt}}$
      & 高噪声区间稳定，DDPM 默认 \\
    $\bx_0$-参数化 & $\bx_0$
      & $\beps=\dfrac{\bx_t-\sqrt{\abt}\,\bx_0}{\sqrt{1-\abt}}$
      & 低噪声区间清晰，高噪声易发散 \\
    $v$-参数化 & $\bm{v}=\sqrt{\abt}\,\beps-\sqrt{1-\abt}\,\bx_0$
      & $\beps=\sqrt{\abt}\,\bm{v}+\sqrt{1-\abt}\,\bx_t$
      & 全区间数值稳定，训练超参更宽容 \\
    \bottomrule
  \end{tabular}
\end{table}

三种参数化在数学上完全等价、可以逐点互相换算，但经验风险关于 $\theta$ 的\textbf{加权方式}不同：
$\beps$-参数化对大 $t$（高噪声）的误差更敏感，$\bx_0$-参数化对小 $t$ 更敏感，
$v$-参数化相当于两端折中，因此在 $T$ 很大或精度要求较高时常常更省事\cite{salimans2022progressive}。

\section{训练与采样算法}

\subsection{训练}

式~\eqref{eq:l-simple} 已经把所有东西都准备好了，训练过程也就没有任何魔法：

\begin{algorithm}[htbp]
  \caption{扩散模型的训练（对应式~\eqref{eq:l-simple}）}
  \label{alg:train}
  \begin{algorithmic}[1]
    \Require 训练集 $\{\bx_0^{(n)}\}\sim p_{\text{data}}$，噪声调度 $\{\beta_t\}_{t=1}^{T}$，网络 $\beps_\theta$
    \Repeat
      \State 采样一个干净样本 $\bx_0\sim p_{\text{data}}$
      \State 采样一个噪声等级 $t\sim\mathcal{U}\{1,\dots,T\}$
      \State 采样一撮噪声 $\beps\sim\Ncal(\bzero,\bI)$
      \State 由闭式解合成含噪样本
             $\bx_t=\sqrt{\abt}\,\bx_0+\sqrt{1-\abt}\,\beps$
      \State 取梯度下降步，梯度为
             $\nabla_\theta\bigl\|\beps-\beps_\theta(\bx_t,t)\bigr\|^2$
    \Until{收敛}
  \end{algorithmic}
\end{algorithm}

注意算法~\ref{alg:train} 的一个特点：\emph{每一轮只加一次噪}。
因为有了闭式解，$t$ 可以任意跳，不需要按顺序走完整条链，
所以每一步的计算代价与 $T$ 无关。这是扩散模型训练成本其实并不高的原因。

\subsection{采样}

采样则是反向走完 $T$ 步：

\begin{algorithm}[htbp]
  \caption{扩散模型的采样（DDPM，$T$ 步）}
  \label{alg:sample}
  \begin{algorithmic}[1]
    \Require 训练好的网络 $\beps_\theta$，噪声调度 $\{\beta_t\}$，反向方差 $\sigma_t^2$
    \Ensure 一个生成样本 $\bx_0$
    \State 采样起点 $\bx_T\sim\Ncal(\bzero,\bI)$
    \For{$t=T,T-1,\dots,1$}
      \If{$t>1$}
        \State 采样 $\bz\sim\Ncal(\bzero,\bI)$
      \Else
        \State $\bz\gets\bzero$
      \EndIf
      \State $\bx_{t-1}\gets\dfrac{1}{\sqrt{\alpha_t}}
             \Bigl(\bx_t-\dfrac{\beta_t}{\sqrt{1-\abt}}\,\beps_\theta(\bx_t,t)\Bigr)
             +\sigma_t\,\bz$
    \EndFor
    \State \textbf{return} $\bx_0$
  \end{algorithmic}
\end{algorithm}

算法~\ref{alg:sample} 的结构与算法~\ref{alg:train} 高度对称：训练是“把已知的噪声加进去”再让网络
把它猜出来，采样是“把预测出来的噪声减掉”。把两段伪代码并排看一遍，
基本就理解扩散模型的全部机制了。

\subsection{反向方差 $\sigma_t^2$ 的取法：随机采样与确定性采样}

DDPM 论文给出两种取法\cite{ho2020ddpm}：
\begin{equation}
  \sigma_t^2=\beta_t
  \qquad\text{或}\qquad
  \sigma_t^2=\tilde\beta_t=\frac{1-\bar\alpha_{t-1}}{1-\abt}\,\beta_t .
  \label{eq:sigma-choices}
\end{equation}
前者在 $\abt$ 变化平缓时与后者接近；后者在 $t=1$ 时恰好为 $0$，意味着最后一步不加噪声。
$\sigma_t$ 实际上是一个“随机性旋钮”：取正数，采样是随机的；取 $\sigma_t=0$，
反向链就退化成一条\emph{确定性}的映射，同一个 $\bx_T$ 永远产生同一个输出。

\begin{remark}[DDIM：把 $1000$ 步压到 $50$ 步]
DDIM 从另一个角度出发\cite{song2021ddim}：它构造了一个\emph{非马尔可夫}的前向过程，
其边缘分布 $q(\bx_t\mid\bx_0)$ 与 DDPM 完全相同，因此可以复用同一个训练好的网络，
但采样式改写为“先估计 $\bx_0$、再重新构造 $\bx_{t-1}$”：
\begin{equation}
  \hat{\bx}_0=\frac{\bx_t-\sqrt{1-\abt}\,\beps_\theta(\bx_t,t)}{\sqrt{\abt}},
  \qquad
  \bx_{t-1}=\sqrt{\bar\alpha_{t-1}}\,\hat{\bx}_0+\sqrt{1-\bar\alpha_{t-1}}\,\beps_\theta(\bx_t,t).
  \label{eq:ddim}
\end{equation}
因为前向过程不再要求马尔可夫性，$\{\bx_t\}$ 可以只在任意\emph{子序列}
$\tau_1<\cdots<\tau_S$ 上取值。于是 $1000$ 步可以被压缩到 $20\sim50$ 步，
且\emph{同一个模型可以随意改变采样步数而无需重新训练}。
式~\eqref{eq:ddim} 与“取 $\sigma_t=0$”在形态上很接近（都没有再注入噪声），
但 DDIM 的跳步能力来自它的非马尔可夫结构，这一点是单纯的 $\sigma_t=0$ 做不到的。
\end{remark}

\subsection{为什么每一步要采样，而不是直接取均值？}

既然网络已经给出了均值 $\bmu_\theta$，每一步\emph{直接取均值}不是更省事吗？这是学习扩散模型时
很自然的一个疑问，它的答案有两层。

\textbf{第一层是概率上的。} $p_\theta(\bx_{t-1}\mid\bx_t)=\Ncal(\bmu_\theta,\sigma_t^2\bI)$
本身是一个\emph{分布}，均值只是它的中心。要“从分布中生成样本”，就必须从分布里采样；
若每一步都取均值，整条反向链变成一个确定性函数 $\bx_0=f(\bx_T)$，
同一个起点永远给出同一张图，众数唯一，样本多样性随之丧失。

\textbf{第二层是生成任务的实践经验。} 在序列生成（文本、语音）中，“每一步都取概率最大的那个词”
会得到重复、单调甚至崩塌的结果，这一现象早在语言模型采样中就有系统讨论
（\emph{The Curious Case of Neural Text Degeneration}\cite{holtzman2020curious}）；
语音合成中也观察到类似现象：完全去掉采样（例如去掉 dropout 带来的随机性）
会让韵律变得机械、不自然\cite{lee-ddpm}。

反过来说，这两点也说明\textbf{均值已经携带了绝大部分信息}。把 $\sigma_t$ 设成 $0$，
生成的图像依然可以接受、甚至更锐利，说明噪声更多影响的是多样性与细节的随机性，
而不是图像的主体结构。所以 $\sigma_t$ 更像是一个“多样性 $\leftrightarrow$ 确定性”的旋钮，
而不是一个非此即彼的选择。

\subsection{一个最小可运行的实现}

把上面两段伪代码直译成 PyTorch，核心不到 40 行（完整可运行版本见 \texttt{code/ddpm\_minimal.py}）：

\begin{lstlisting}[caption={DDPM 的最小实现（训练与采样）}, label={lst:ddpm}]
import torch

T = 1000
beta  = torch.linspace(1e-4, 2e-2, T)      # linear noise schedule
alpha = 1.0 - beta
abar  = torch.cumprod(alpha, dim=0)        # cumulative product of alpha

def q_sample(x0, t, noise):
    """Forward closed form: x_t = sqrt(abar_t) x0 + sqrt(1-abar_t) noise"""
    a = abar[t].sqrt().view(-1, *[1] * (x0.dim() - 1))
    s = (1.0 - abar[t]).sqrt().view(-1, *[1] * (x0.dim() - 1))
    return a * x0 + s * noise

def train_step(model, x0, opt):
    t     = torch.randint(0, T, (x0.shape[0],), device=x0.device)
    noise = torch.randn_like(x0)
    xt    = q_sample(x0, t, noise)              # one shot, O(1)
    loss  = ((model(xt, t) - noise) ** 2).mean() # predict the noise
    opt.zero_grad(); loss.backward(); opt.step()
    return loss.item()

@torch.no_grad()
def sample(model, shape):
    x = torch.randn(shape)                       # x_T ~ N(0, I)
    for ti in reversed(range(T)):
        t     = torch.full((shape[0],), ti, dtype=torch.long)
        eps   = model(x, t)
        mean  = (x - beta[ti] / (1.0 - abar[ti]).sqrt() * eps) / alpha[ti].sqrt()
        if ti > 0:
            x = mean + beta[ti].sqrt() * torch.randn_like(x)   # sigma_t^2 = beta_t
        else:
            x = mean
    return x
\end{lstlisting}

\section{延伸与应用}

\subsection{连续化的视角：随机微分方程}

把 $T$ 取得很大、把 $\beta_t$ 取得很小的极限下，前向过程就是一个\emph{随机微分方程}（SDE）
\begin{equation}
  \mathrm{d}\bx = -\tfrac{1}{2}\beta(t)\,\bx\,\mathrm{d}t+\sqrt{\beta(t)}\,\mathrm{d}\bm{w},
  \label{eq:forward-sde}
\end{equation}
其中 $\bm{w}$ 是标准布朗运动。它的反向时间 SDE 为
\begin{equation}
  \mathrm{d}\bx=\Bigl[-\tfrac{1}{2}\beta(t)\,\bx-\beta(t)\,\nabla_{\bx}\log q_t(\bx)\Bigr]\mathrm{d}t
  +\sqrt{\beta(t)}\,\mathrm{d}\bar{\bm{w}},
  \label{eq:reverse-sde}
\end{equation}
即只需要知道每个时刻的得分函数 $\nabla_{\bx}\log q_t(\bx)$；而由上一节的结论，
得分恰好由噪声预测网络给出。Song 等人指出，同一族前向 SDE 都对应一个\emph{概率流 ODE}
\begin{equation}
  \mathrm{d}\bx=\Bigl[-\tfrac{1}{2}\beta(t)\,\bx-\tfrac{1}{2}\beta(t)\,\nabla_{\bx}\log q_t(\bx)\Bigr]\mathrm{d}t,
  \label{eq:pf-ode}
\end{equation}
沿着它积分能得到完全确定性的采样轨迹，且与 SDE 采样共享同一组边缘分布\cite{song2021sde}。
这从理论上解释了 DDIM 的确定性采样与跳步能力：它正是概率流 ODE 的一种数值离散。

\subsection{语音与文本}

扩散模型最初在图像上取得突破，随后被搬到其他模态，而不同模态的难点恰好揭示了它的边界。

\textbf{语音}几乎没有障碍，因为语音波形本身是连续的高维实值信号，
与图像的设定同构。WaveGrad、DiffWave 等模型把扩散用于声码器与声学模型，
用几十步采样即可合成高质量语音\cite{chen2021wavegrad}。

\textbf{文本}则要麻烦得多，原因是文本是\emph{离散}的：
标准高斯噪声无法与 token 的 one-hot 表示直接相加，而前向过程一旦不再可加高斯，
闭式解与整套推导都会失效。实践中形成了三条路：

\begin{itemize}
  \item \textbf{在隐空间上加噪}：先用一个连续的自编码器把文本（或图像）压到连续的隐空间，
        再在隐空间上做标准扩散。Latent Diffusion 在图像上非常成功，
        DiffuSeq 等把这一思路搬到文本，可同时生成一对相关的句子\cite{rombach2022ldm}；
  \item \textbf{换掉高斯噪声}：不用高斯噪声，而用“编辑”（插入、删除、替换）这类\emph{离散}的
        扰动来定义前向过程，DiffusER 是这一路线的代表；
  \item \textbf{用掩码代替加噪}：把“加噪—去噪”换成“加掩码—预测”，
        即 Mask-Predict 一类的非自回归生成方法，可以并行地一次填空若干位置\cite{ghazvininejad2019mask}。
\end{itemize}

\begin{remark}[一个统一的理解]
如果把“扩散”理解为\emph{把一次生成拆成若干次微小的、逐步确定化的生成}，
那么这些变体其实是同一件事在不同数据模态上的实现：
连续数据用高斯噪声，离散数据用掩码或编辑操作，
维度太高（例如高分辨率图像）就先压到隐空间再扩散。
\end{remark}

\subsection{条件生成与工程实践}

要控制生成结果（“一只在奔跑的狗”），只需把条件 $\bm{c}$ 交给网络：
$\beps_\theta(\bx_t,t)\to\beps_\theta(\bx_t,t,\bm{c})$，训练时把 $\bm{c}$（文本嵌入、类别标签）
与时间步 $t$ 一起注入。更细的控制则通过\emph{引导}（guidance）实现：

\begin{itemize}
  \item \textbf{分类器引导}（classifier guidance）：额外训一个分类器 $p_\phi(\bm{c}\mid\bx_t)$，
        在采样时把它的梯度加到均值上，$\tilde{\bmu}=\bmu_\theta+\sigma_t^2\nabla_{\bx_t}\log p_\phi(\bm{c}\mid\bx_t)$；
  \item \textbf{无分类器引导}（classifier-free guidance）：不训分类器，
        而是同时训练“有条件”与“无条件”两种预测，采样时做外推
        $\tilde{\beps}=(1+w)\,\beps_\theta(\bx_t,t,\bm{c})-w\,\beps_\theta(\bx_t,t,\varnothing)$，
        $w$ 越大越贴合条件、但多样性越低\cite{ho2022cfg}。
\end{itemize}

最后是工程上最常被提到的三个议题：

\begin{itemize}
  \item \textbf{采样太慢}：$T$ 次网络前向是主要成本。加速思路有三类——
        改采样器（DDIM、DPM-Solver 等把步数压到几十步甚至十几步）、
        蒸馏（Progressive Distillation、一致性模型，把多步压成 1～4 步）、
        在隐空间扩散（把每次前向的分辨率降下来）；
  \item \textbf{损失加权}：式~\eqref{eq:Lt-weighted} 中的 $w_t$ 丢掉之后虽然效果好，
        但“如何加权”本身是一个可以调的设计空间，$\beps$/$v$ 参数化与 Min-SNR 加权都属于这一类；
  \item \textbf{训练数据与算力}：扩散模型的训练是稳定的（没有 GAN 那样的对抗失衡），
        代价转移到数据规模与采样时间上，这也是它在这几年迅速工业化的原因之一。
\end{itemize}

\section{小结}

把整篇内容压成一条线，扩散模型的逻辑其实非常顺：

\begin{enumerate}
  \item \textbf{目标}：生成模型想学 $p_{\text{data}}(\bx)$，最大似然等价于最小化
        $\Dkl(p_{\text{data}}\|p_\theta)$，但 $p_\theta(\bx)$ 里的积分算不出来；
  \item \textbf{方案}：把“一步生成”改成“$T$ 步微小的去噪”。用一条\emph{固定、无参数}的前向链
        把数据逐步溶成噪声（相当于 VAE 的编码器），再学一条反向链把它还原（解码器）；
  \item \textbf{前向闭式解}：$q(\bx_t\mid\bx_0)=\Ncal(\sqrt{\abt}\bx_0,(1-\abt)\bI)$
        让 $\bx_t$ 可以\emph{一步}造出来，训练代价与 $T$ 无关；
  \item \textbf{下界}：由 Jensen 不等式得 ELBO，展开成 $L_T+\sum_{t\ge2}L_{t-1}+L_0$。
        其中 $L_T$ 是常数，$L_{t-1}$ 是主项，而真实后验
        $q(\bx_{t-1}\mid\bx_t,\bx_0)$ 有闭式解；
  \item \textbf{化简}：两个方差相同的正态分布之间的 KL 退化成均值平方误差，
        再把均值换成“预测噪声”的写法，就得到
        $\mathcal{L}_{\text{simple}}=\E\|\beps-\beps_\theta(\bx_t,t)\|^2$；
  \item \textbf{采样}：从 $\Ncal(\bzero,\bI)$ 出发反向走 $T$ 步，每步减掉预测的噪声、
        再按 $\sigma_t$ 补一点新噪声。$\sigma_t$ 控制随机性；
        $\sigma_t\to0$ 得到确定性采样，DDIM 由此把步数压到几十步；
  \item \textbf{联系}：预测噪声等价于估计得分函数，于是扩散模型可以写成反向时间 SDE 与概率流 ODE，
        也解释了它与 VAE、得分匹配之间的关系。
\end{enumerate}

\begin{keybox}[一句话]
\textbf{扩散模型 = 用一条固定的加噪链，把一个不可解的积分问题，
换成了一个可以解析写出下界的、逐层的去噪回归问题。}
它的所有“魔法”都来自前向过程的闭式解 $q(\bx_t\mid\bx_0)$：
正因为 $\bx_t$ 可以一步造出来，训练才会既简单又便宜；
也正因为真实后验可以解析写出，变分下界才能化简成“看图猜噪声”。
\end{keybox}

表~\ref{tab:cheatsheet} 是几个核心公式的速查表，动手实现时可以据此逐行对照。

\begin{table}[htbp]
  \centering
  \caption{核心公式速查}
  \label{tab:cheatsheet}
  \small
  \renewcommand{\arraystretch}{1.85}
  \setlength{\tabcolsep}{5pt}
  \begin{tabular}{@{}p{3.3cm} p{10.2cm}@{}}
    \toprule
    环节 & 公式 \\
    \midrule
    前向单步     & $\bx_t=\sqrt{\alpha_t}\,\bx_{t-1}+\sqrt{\beta_t}\,\beps$ \\
    前向闭式     & $\bx_t=\sqrt{\abt}\,\bx_0+\sqrt{1-\abt}\,\beps$ \\
    后验均值     & $\tilde{\bmu}_t=\dfrac{\sqrt{\bar\alpha_{t-1}}\,\beta_t\,\bx_0+\sqrt{\alpha_t}\,(1-\bar\alpha_{t-1})\,\bx_t}{1-\abt}$ \\
    训练目标     & $\mathcal{L}_{\text{simple}}=\E\bigl\|\beps-\beps_\theta(\bx_t,t)\bigr\|^2$ \\
    采样单步     & $\bx_{t-1}=\dfrac{1}{\sqrt{\alpha_t}}\Bigl(\bx_t-\dfrac{\beta_t}{\sqrt{1-\abt}}\,\beps_\theta(\bx_t,t)\Bigr)+\sigma_t\bz$ \\
    从噪声还原数据 & $\hat{\bx}_0=\dfrac{\bx_t-\sqrt{1-\abt}\,\beps_\theta(\bx_t,t)}{\sqrt{\abt}}$ \\
    噪声与得分   & $\nabla_{\bx_t}\log q(\bx_t)=-\,\dfrac{\beps_\theta(\bx_t,t)}{\sqrt{1-\abt}}$ \\
    \bottomrule
  \end{tabular}
\end{table}


\appendix
\section{两个高斯分布的常用结论}

\subsection{独立高斯之和仍是高斯}
\label{app:gauss-sum}

\begin{proposition}
设 $\beps_1\sim\Ncal(\bzero,\sigma_1^2\bI)$、$\beps_2\sim\Ncal(\bzero,\sigma_2^2\bI)$ 相互独立，
$a,b\in\mathbb{R}$，则
\begin{equation}
  a\,\beps_1+b\,\beps_2\ \sim\ \Ncal\bigl(\bzero,\ (a^2\sigma_1^2+b^2\sigma_2^2)\bI\bigr).
  \label{eq:gauss-sum}
\end{equation}
\end{proposition}

\begin{pf}
设 $\bm{t}\in\mathbb{R}^d$，用特征函数。由于 $\beps_1,\beps_2$ 独立，
\begin{equation}
  \E\bigl[e^{\mathrm{i}\bm{t}^{\top}(a\beps_1+b\beps_2)}\bigr]
  =\E\bigl[e^{\mathrm{i}a\bm{t}^{\top}\beps_1}\bigr]\cdot\E\bigl[e^{\mathrm{i}b\bm{t}^{\top}\beps_2}\bigr]
  =e^{-\frac{1}{2}a^2\sigma_1^2\|\bm{t}\|^2}\cdot e^{-\frac{1}{2}b^2\sigma_2^2\|\bm{t}\|^2}
  =e^{-\frac{1}{2}(a^2\sigma_1^2+b^2\sigma_2^2)\|\bm{t}\|^2},
\end{equation}
这正是零均值高斯 $\Ncal(\bzero,(a^2\sigma_1^2+b^2\sigma_2^2)\bI)$ 的特征函数，
由特征函数与分布的一一对应即得结论。
\end{pf}

这个结论在正文中被用了两次：一次是在前向闭式解的归纳法里
（式~\eqref{eq:forward-induction}），把两步噪声合并成一步；
另一次是在推导 $\tilde{\bmu}_t$ 时把“关于 $\bx_{t-1}$ 的二次型”配方成高斯。

\subsection{两个高斯之间的 KL 散度}
\label{app:kl-gauss}

\begin{proposition}
设 $\bmu_1,\bmu_2\in\mathbb{R}^d$，$\bm{\Sigma}_1,\bm{\Sigma}_2$ 为 $d\times d$ 正定矩阵，则
\begin{equation}
  \begin{aligned}
  \Dkl\bigl(\Ncal(\bmu_1,\bm{\Sigma}_1)\,\big\|\,\Ncal(\bmu_2,\bm{\Sigma}_2)\bigr)
  =\frac{1}{2}\Bigl[&\operatorname{tr}\bigl(\bm{\Sigma}_2^{-1}\bm{\Sigma}_1\bigr)
   +(\bmu_2-\bmu_1)^{\top}\bm{\Sigma}_2^{-1}(\bmu_2-\bmu_1) \\
   &-d+\ln\frac{\det\bm{\Sigma}_2}{\det\bm{\Sigma}_1}\Bigr].
  \end{aligned}
  \label{eq:kl-gauss}
\end{equation}
特别地，当两者协方差相同（$\bm{\Sigma}_1=\bm{\Sigma}_2=\sigma^2\bI$）时，
只有均值不同，公式退化成
\begin{equation}
  \Dkl\bigl(\Ncal(\bmu_1,\sigma^2\bI)\,\big\|\,\Ncal(\bmu_2,\sigma^2\bI)\bigr)
  =\frac{1}{2\sigma^2}\bigl\|\bmu_1-\bmu_2\bigr\|^2 .
  \label{eq:kl-gauss-equal-var}
\end{equation}
\end{proposition}

式~\eqref{eq:kl-gauss-equal-var} 就是正文式~\eqref{eq:Lt-kl} 的来源：
它把“两个分布有多远”这件事，换成了“两个均值差多少”这件更简单的事，
而后者正好可以用平方误差来优化。

\subsection{ELBO 三项分解的完整展开}
\label{app:elbo}

从式~\eqref{eq:ddpm-elbo} 出发，把两个链式乘积代入：
\begin{equation}
  \begin{aligned}
    \mathcal{L}(\bx_0)
    &=\E_q\Bigl[\log p(\bx_T)+\sum_{t=1}^{T}\log p_\theta(\bx_{t-1}\mid\bx_t)
      -\sum_{t=1}^{T}\log q(\bx_t\mid\bx_{t-1})\Bigr] \\
    &=\E_q\Bigl[\log p(\bx_T)+\log p_\theta(\bx_0\mid\bx_1) \\
    &\qquad\qquad+\sum_{t=2}^{T}\log p_\theta(\bx_{t-1}\mid\bx_t)
      -\sum_{t=2}^{T}\log q(\bx_t\mid\bx_{t-1})-\log q(\bx_1\mid\bx_0)\Bigr].
  \end{aligned}
  \label{eq:app-elbo-1}
\end{equation}
对 $t\ge2$ 的每一对相邻项用贝叶斯公式~\eqref{eq:bayes-trick}：
\begin{equation}
  \begin{aligned}
    \sum_{t=2}^{T}\log q(\bx_t\mid\bx_{t-1})
    &=\sum_{t=2}^{T}\log q(\bx_{t-1}\mid\bx_t,\bx_0) \\
    &\quad+\underbrace{\sum_{t=2}^{T}\bigl[\log q(\bx_t\mid\bx_0)-\log q(\bx_{t-1}\mid\bx_0)\bigr]}_{\text{相消}},
  \end{aligned}
  \label{eq:app-elbo-2}
\end{equation}
后一个求和是相邻两项的差，只剩首尾：$\log q(\bx_T\mid\bx_0)-\log q(\bx_1\mid\bx_0)$。
代回式~\eqref{eq:app-elbo-1}，其中 $-\log q(\bx_1\mid\bx_0)$ 恰好被抵消，于是
\begin{equation}
  \mathcal{L}(\bx_0)
  =\E_q\Bigl[\log p_\theta(\bx_0\mid\bx_1)
   -\sum_{t=2}^{T}\log\frac{q(\bx_{t-1}\mid\bx_t,\bx_0)}{p_\theta(\bx_{t-1}\mid\bx_t)}
   -\log\frac{q(\bx_T\mid\bx_0)}{p(\bx_T)}\Bigr].
  \label{eq:app-elbo-3}
\end{equation}
现在逐项看：对固定的 $t$，第二项里除 $\bx_t$ 之外的变量都能积掉，期望退化为
$\E_{q(\bx_t\mid\bx_0)}\bigl[\Dkl(q(\bx_{t-1}\mid\bx_t,\bx_0)\|p_\theta(\bx_{t-1}\mid\bx_t))\bigr]$；
第三项就是 $\Dkl(q(\bx_T\mid\bx_0)\|p(\bx_T))$；第一项是 $-\log p_\theta(\bx_0\mid\bx_1)$ 的期望。于是
\begin{equation}
  \begin{aligned}
    \mathcal{L}(\bx_0)
    &=-\underbrace{\Dkl\bigl(q(\bx_T\mid\bx_0)\|p(\bx_T)\bigr)}_{L_T} \\
    &\quad-\sum_{t=2}^{T}\underbrace{\E_{q(\bx_t\mid\bx_0)}
      \bigl[\Dkl\bigl(q(\bx_{t-1}\mid\bx_t,\bx_0)\|p_\theta(\bx_{t-1}\mid\bx_t)\bigr)\bigr]}_{L_{t-1}} \\
    &\quad-\underbrace{\E\bigl[-\log p_\theta(\bx_0\mid\bx_1)\bigr]}_{L_0},
  \end{aligned}
  \label{eq:app-elbo-4}
\end{equation}
即 $\mathcal{L}(\bx_0)=-(L_T+\sum_{t=2}^{T}L_{t-1}+L_0)$，与正文式~\eqref{eq:elbo-three} 一致。


\printbibliography[title={参考文献}]

\end{document}
